\section{Welfare}\label{sec:welfare}

The operator microfoundation provides a direct welfare interpretation of the same access feedback that shapes price competition. Let $\bar v$ denote the gross value of a trip for each background passenger. Shopper consumer surplus is
\begin{align}
CS_S
={}&\int_0^x\left[v-p_L-tz^2-\frac{w}{f_L}\right]dz
+\int_x^1\left[v-p_R-t(1-z)^2-\frac{w}{f_R}\right]dz\\
={}&v-p_Lx-p_R(1-x)
-\frac{t}{3}\{x^3+(1-x)^3\}
-A\left\{\frac{x}{s_q}+\frac{1-x}{1-s_q}\right\}.
\label{eq:shopper_cs}
\end{align}
Background-passenger surplus is
\begin{equation}
CS_B=M\bar v-A\frac{M}{s_q},
\end{equation}
and retail producer surplus is
\begin{equation}
PS=(p_L-c)x+(p_R-c)(1-x).
\end{equation}
Adding these terms cancels retail prices exactly. Up to constants, total welfare is therefore the negative of real spatial and waiting costs:
\begin{equation}
TW=v+M\bar v-c-C_q(x),
\label{eq:TW}
\end{equation}
where
\begin{equation}
C_q(x)=
\frac{t}{3}\{x^3+(1-x)^3\}
+A\left\{\frac{M+x}{s_q(x,M)}+\frac{1-x}{1-s_q(x,M)}\right\}.
\label{eq:real_cost}
\end{equation}
Any fixed cost of supplying the resource $F$ is common across the comparisons considered here and can be added as a constant.

\subsection{The operator envelope identity}

When the service floor is slack, minimized aggregate waiting cost is
\begin{equation}
J(x)=A\left(\sqrt{M+x}+\sqrt{1-x}\right)^2.
\label{eq:J}
\end{equation}
The operator envelope theorem gives
\begin{equation}
J'(x)=A H(x,M).
\label{eq:envelope}
\end{equation}
Consequently,
\begin{equation}
C'(x)=t(2x-1)+AH(x,M),
\label{eq:Cprime}
\end{equation}
and the marginal shopper condition \eqref{eq:Gq} can be written on the slack branch as
\begin{equation}
p_L-p_R+C'(x)=0.
\label{eq:demand_welfare_identity}
\end{equation}

\begin{proposition}[Operator-envelope welfare identity]\label{prop:envelope}
When the minimum service constraint is slack, the directional waiting-cost difference entering shopper choice equals the derivative of minimized aggregate waiting cost. Hence the access term that determines retail demand also has an exact real-cost counterpart in welfare.
\end{proposition}

This identity is the main reason to retain the transport operator rather than replace it with an arbitrary market-share externality. A reduced-form externality can reproduce local demand curvature, but it does not by itself determine the associated resource or waiting-cost accounting.

\subsection{Private and social allocation}

On a stable slack branch, $C''(x)=g(x)>0$. Combining \eqref{eq:demand_welfare_identity} with the price first-order conditions \eqref{eq:price_focs} gives
\begin{equation}
C'(x^{NE})=-(2x^{NE}-1)C''(x^{NE}).
\label{eq:welfare_wedge}
\end{equation}

\begin{proposition}[Private--social share wedge]\label{prop:welfare_wedge}
At a stable interior price equilibrium with a slack service floor, if $x^{NE}>1/2$ then $C'(x^{NE})<0$, so a marginal reallocation of shoppers toward $L$ lowers real social cost. If $x^{NE}<1/2$, the symmetric statement holds for $R$. The local wedge vanishes at $x^{NE}=1/2$.
\end{proposition}

The decentralized price game therefore dampens the socially valuable concentration toward the direction that already carries greater total demand. The economic force is the service-frequency scale economy generated by sharing a fixed fleet across shoppers and background passengers.

At the exact witness, $x^{NE}=0.575$ and
\begin{equation}
C'(x^{NE})=-\frac{282351}{1841170}\approx-0.153354.
\end{equation}
Under the same service floor $q=1/3$, the right-direction floor binds for sufficiently high $x$. In that region the unique real-cost minimizer is
\begin{equation}
x^{SB}(q)=\frac12+\frac{A(1-2q)}{2tq(1-q)}.
\label{eq:second_best_general}
\end{equation}
For the witness,
\begin{equation}
x^{SB}\approx0.702378.
\end{equation}
The associated reduction in real cost is approximately $0.011378$, about $1.22\%$ of decentralized real cost under the witness normalization.

\input{tables/welfare_summary}

The appropriate welfare comparison in this paper is the decentralized equilibrium versus the \emph{same-floor} second best. If the service obligation is removed from the planner problem at the witness, the utilitarian first best places all shoppers on the $L$ side and shuts down service in the other direction. We do not interpret that corner as a policy recommendation. Rather, it shows that the service floor embodies a coverage or service-availability objective that is exogenous to the utilitarian objective in \eqref{eq:TW}. Accordingly, the model does not derive an optimal coverage standard or attach an equity value to minimum service.
